\section{Introduction}
\label{sec:intro}
Predictive maintenance models use machine telemetry to estimate the Remaining Useful Life (RUL) of equipment. Telemetry can be operationally sensitive, while individual sites may have limited examples of degradation and failure. Federated learning (FL) lets clients train on local data and share model updates with a coordinator, but exchanging model updates is not itself a formal privacy guarantee~\cite{mcmahan2017fedavg,abadi2016dpsgd}. Practical RUL work must therefore distinguish data locality from protections such as differential privacy and secure aggregation.

The C-MAPSS benchmark is well established for RUL prediction, and centralized CNN, recurrent, and attention-based approaches report strong results across its four subsets~\cite{li2018cnn,zheng2017lstm,wang2024adamsamtcn,liu2025moeformer}. Federated RUL studies already include C-MAPSS experiments, FedAvg/FedProx, personalization layers, and pruning~\cite{rosero2020federated,soderkvist2024collaborative,chen2023pruning}, while recent N-CMAPSS and C-MAPSS studies extend this work to airline collaboration and controlled heterogeneous benchmarks~\cite{landau2025aircraftfl,sorrenti2026fedcmapss}.

RUL evaluation also needs metrics that reflect asymmetric maintenance costs. A late prediction ($\hat{y} > y$) may defer service beyond the actual failure point, while an early prediction ($\hat{y} < y$) may consume useful life. The NASA/PHM08 scoring function penalizes late errors more heavily than early errors~\cite{saxena2008damage,saxena2008metrics}; reporting RMSE alone can therefore hide relevant differences. We evaluate two risk-aware objectives without treating them as the same experiment: the paired FD001--FD004 suite includes a capped exponential NASA-score objective, whereas a separate FD001/FD002 study compares MSE with \texttt{AsymmetricRULLoss}, a $C^1$-smooth asymmetric quadratic surrogate. Both report NASA score and RMSE, and the latter also reports late-prediction rate.

A low RUL estimate alone does not tell a maintainer what to inspect. The edge node already holds the raw telemetry and the model, so it can also explain the estimate locally: which engine module (fan, low- or high-pressure compressor, hot section) and which sensors are associated with the drop. Doing this on the node keeps raw sensor values off the network. C-MAPSS, however, publishes unit-level RUL only, with no per-unit component fault labels~\cite{saxena2008damage}. Any such explanation can therefore be checked for physical consistency and unsupervised separability, but not scored as supervised component-diagnosis accuracy.

This paper follows one edge node through three steps: it learns with other sites, it trains against asymmetric risk, and it explains its estimate locally. It asks: (1) how do centralized, local-only and FedAvg models compare under matched preprocessing and seeds, and how do client partition and aggregation method affect global and per-client accuracy; (2) how does asymmetric-risk optimization alter symmetric accuracy and late-prediction risk; and (3) can a local attribution and drift module point to physically plausible subsystems at an acceptable per-window cost, given that no component labels exist? The contributions are:
\begin{itemize}
\item a paired FD001--FD004 benchmark of centralized, local-only, FedAvg, FedProx and FedAdam training, with a capped exponential score objective and a compact CNN-LSTM ablation;
\item a separate controlled five-seed evaluation of quadratic \texttt{AsymmetricRULLoss} across local-only, FedAvg and centralized setups on FD001 and FD002;
\item a local degradation attribution and healthy-baseline drift module that maps the model's 14 sensor channels to four subsystems, excludes operating settings from fault responsibility, and makes no network calls; and an offline dashboard that displays its output;
\item a five-seed FD001--FD004 validation of that module against sensor-count chance rates, unsupervised two-mode separation and single-thread latency.
\end{itemize}
An edge profiler is also provided, but this manuscript reports only its committed Mac smoke profile and the diagnostics latency on an Apple Silicon host. It does not claim Raspberry Pi efficiency, encrypted transport, formal privacy or supervised component-diagnosis accuracy.

The remainder is organised as follows: Section~\ref{sec:related} reviews related work, Section~\ref{sec:arch} the architecture, Section~\ref{sec:method} the method, Section~\ref{sec:setup} the setup, Section~\ref{sec:results} the results, Section~\ref{sec:discussion} limitations, and Section~\ref{sec:conclusion} concludes.
